\begin{titlepage}
\thispagestyle{empty}
\begin{tikzpicture}[remember picture, overlay]
  \fill[inkgray] (current page.north west) rectangle (current page.south east);
  \fill[mossgreen!24] ([xshift=-0.7cm,yshift=-2.0cm]current page.north east) circle (6.4cm);
  \fill[white!8] ([xshift=1.1cm,yshift=2.0cm]current page.south west) circle (5.2cm);

  \draw[warmgold!85, line width=1.4pt]
    ([xshift=1.1cm,yshift=-1.1cm]current page.north west)
    rectangle
    ([xshift=-1.1cm,yshift=1.1cm]current page.south east);

  \node[
    anchor=north west,
    fill=warmgold,
    text=inkgray,
    font=\bfseries\large,
    rounded corners=4pt,
    inner xsep=10pt,
    inner ysep=6pt
  ] at ([xshift=1.6cm,yshift=-1.5cm]current page.north west) {QUYỂN XVII};

  \node[
    anchor=north,
    text=white,
    font=\large\itshape
  ] at ([yshift=-1.7cm]current page.north) {Truyện Tích Cảm Hứng Song Ngữ};

  \node[
    anchor=west,
    text=white,
    font=\fontsize{21}{25}\selectfont\bfseries
  ] at ([xshift=1.8cm,yshift=4.3cm]current page.center) {NHỮNG ĐIỀU};

  \node[
    anchor=west,
    text=warmgold,
    font=\fontsize{21}{25}\selectfont\bfseries
  ] at ([xshift=1.8cm,yshift=2.4cm]current page.center) {HỌC SINH NÊN};

  \node[
    anchor=west,
    text=white,
    font=\fontsize{21}{25}\selectfont\bfseries
  ] at ([xshift=1.8cm,yshift=0.5cm]current page.center) {HIỂU SỚM};

  \draw[warmgold, line width=2pt]
    ([xshift=1.8cm,yshift=-1.0cm]current page.center) --
    ([xshift=8.6cm,yshift=-1.0cm]current page.center);

  \node[
    anchor=west,
    text=white,
    font=\large
  ] at ([xshift=1.8cm,yshift=-2.1cm]current page.center) {Một quyển sách gom lại những điều căn bản nhưng rất quan trọng mà học sinh hiểu càng sớm thì càng đỡ vấp về sau};

  \node[
    anchor=north west,
    fill=white,
    text=black,
    rounded corners=8pt,
    text width=11.7cm,
    inner sep=14pt,
    align=left,
    font=\normalsize
  ] at ([xshift=1.9cm,yshift=-14.9cm]current page.north west) {%
    \textbf{Tinh thần quyển này:} nói lại những điều đúng rất sớm, rất nền, để học sinh có một điểm tựa rõ hơn cho việc học, cách nghĩ, tình bạn và đường tương lai.\\[6pt]
    \textbf{Cách dùng:} đọc từng bài ngắn mỗi ngày, suy nghĩ thêm bằng vài câu tiếng Anh đơn giản, và dùng như một cuốn sổ nhỏ để tự nhắc mình.\\[6pt]
    \textbf{Điểm tựa:} gần gũi, ngắn gọn, không giáo điều, nhưng giữ được phần chín chắn và lâu bền.
  };

  \node[
    anchor=south,
    text=warmgold,
    font=\large\itshape,
    align=center
  ] at ([yshift=2.1cm]current page.south) {%
    ``Có những điều hiểu sớm sẽ không làm tuổi học trò nặng hơn,\\
    mà chỉ làm con đường lớn lên bớt mù mờ hơn.''%
  };

  \node[
    anchor=south,
    text=white!85,
    font=\normalsize,
    align=center
  ] at ([yshift=1.1cm]current page.south) {Dành cho học sinh, giáo viên và những ai muốn giữ lại những nền tảng quan trọng nhất của việc học và trưởng thành};
\end{tikzpicture}
\end{titlepage}
\clearpage